\textbf{UE :} S-INFO-015 Projet d'analyse et de conception logicielle\\
\textbf{Année académique :} 2026-2027\\
\textbf{Groupe :} G02 \textbf{Responsable de l'extension :} Fontaine
Alexis (Matricule : 250602)

\hypertarget{extension-gestion-de-cartes}{%
\subsubsection{Extension : gestion de
cartes}\label{extension-gestion-de-cartes}}

\hypertarget{us1-demander-une-carte}{%
\subparagraph{US1 Demander une carte}\label{us1-demander-une-carte}}

En tant que titulaire, je veux pouvoir demander une nouvelle carte à une
institution,\\
afin d'y ascocier un compte.

\textbf{Critères d'acceptation :}

\begin{itemize}
\tightlist
\item
  L'utilisateur doit-être conncecté.
\item
  Chaque carte a un numéro unique, une date d'expiration, un code de
  sécurité et un code PIN.
\item
  La carte reçue est inactive.
\end{itemize}

\hypertarget{us2-lier-un-compte-uxe0-une-carte}{%
\subparagraph{US2 Lier un compte à une
carte}\label{us2-lier-un-compte-uxe0-une-carte}}

En tant que titulaire possédant une carte active, je veux pouvoir y
ascocier un compte, afin de faire des achats avec la carte.

\textbf{Critères d'acceptation :}

\begin{itemize}
\tightlist
\item
  Une carte n'est liée qu'a un seul compte.
\item
  Un compte peut être lié à plusieurs cartes.
\item
  L'utilisateur doit être titulaire ou co-titulaire du compte.
\item
  Il doit être possible de dissocier une carte d'un compte.
\end{itemize}

\hypertarget{us3-visualisation-de-la-carte}{%
\subparagraph{US3 Visualisation de la
carte}\label{us3-visualisation-de-la-carte}}

En tant que titulaire possédant une carte, je veux pouvoir voir les
informations la concernant (date, numéro, code de sécurité/verification)
ainsi que le nombre de crédit restant pour les cartes de crédits,\\
afin de savoir l'argent que je peux encore dépenser.

\textbf{Critères d'acceptation :}

\begin{itemize}
\tightlist
\item
  La visualisation n'est possible que si l'utilisateur est connecté.
\item
  Précise si la carte est active ou non.
\end{itemize}

\hypertarget{us4-se-connecter-avec-une-carte}{%
\subparagraph{US4 Se connecter avec une
carte}\label{us4-se-connecter-avec-une-carte}}

En tant que titulaire possédant une carte active, je veux pouvoir me
connecter à l'application avec mon numéro de carte et un code PIN,\\
afin de faciliter la connection.

\textbf{Critères d'acceptation :}

\begin{itemize}
\tightlist
\item
  Il s'agit d'une fonctionnalité que l'utilsateur peut
  activer/désactiver pour chacune de ses cartes.
\item
  Il doit y avoir une limite maximale d'échecs pour entrer le code PIN.
  Si la limite est dépassée, une notification est envoyée et cette carte
  ne peut plus servir pour se connecter à l'application temps que
  l'utilisateur ne réactive pas cette fonctionnalité.
\end{itemize}

\hypertarget{us5-visualisation-de-lhistorique-de-la-carte}{%
\subparagraph{US5 Visualisation de l'historique de la
carte}\label{us5-visualisation-de-lhistorique-de-la-carte}}

En tant que titulaire possédant une carte, je veux pouvoir visualiser et
exporter l'historique d'une carte de la même façon que je visualise
l'historique d'un compte, afin de stocker les données en local.

\textbf{Critères d'acceptation :}

\begin{itemize}
\tightlist
\item
  La visualisation doit être la même que celle utilisée pour les
  comptes.
\item
  Une carte est considéré comme un produit financier pour la
  visualisation.
\end{itemize}

\hypertarget{us6-expiration-dune-carte}{%
\subparagraph{US6 Expiration d'une
carte}\label{us6-expiration-dune-carte}}

En tant que titulaire possédant une carte active, je veux
automatiquement recevoir une nouvelle carte lorsque l'une d'entre elle
expire, afin de continuer mes achats.

\textbf{Critères d'acceptation :}

\begin{itemize}
\tightlist
\item
  Une carte inactive qui expire, ne déclenche pas la procédure.
\item
  Une fois la date d'expiration dépassée, l'utilisation de la carte est
  impossible.
\item
  Avant l'expiration de la carte, l'utilisateur est averti et reçoit une
  nouvelle carte non active.
\item
  Lors de l'activation de la nouvelle carte, l'ancienne est
  automatiquement désactivée.
\end{itemize}

\hypertarget{us7-bloquer-une-carte}{%
\subparagraph{US7 Bloquer une carte}\label{us7-bloquer-une-carte}}

En tant que titulaire possédant une carte, je veux pouvoir bloquer une
carte et au besoin en demander une nouvelle, afin de me préserver en cas
de perte ou de vol.

\textbf{Critères d'acceptation :}

\begin{itemize}
\tightlist
\item
  Il est possible de bloquer une carte qui n'est pas encore active.
\item
  La demande doit être envoyée à la banque qui bloquera la carte.
\item
  L'utilisateur reçoit à la demande, une nouvelle carte qu'il devra
  activer.
\item
  Une carte bloquée est inactive et ne peut pas être réactivée.
\end{itemize}

\hypertarget{us8-modifier-une-carte}{%
\subparagraph{US8 Modifier une carte}\label{us8-modifier-une-carte}}

En tant que titulaire possédant une carte active, je veux pouvoir
modifier le seuil de crédit accordé, autoriser ou non l'utilisation à
l'étranger et autoriser ou non les transactions qui donneront lieues à
un solde négatif, afin d'utiliser ma carte dans les meilleures
conditions.

\textbf{Critères d'acceptation :}

\begin{itemize}
\tightlist
\item
  Modifier le seuil de crédit doit être validé par la banque.
\item
  Les modifications peuvent être temporaire avec une date de fin
  prédéfinie.
\end{itemize}

\hypertarget{us9-utiliser-la-carte}{%
\subparagraph{US9 Utiliser la carte}\label{us9-utiliser-la-carte}}

En tant que titulaire possédant une carte active, je veux pouvoir
utiliser ma carte pour effectuer des achats,\\
afin de le faire de façon simplifié.

\textbf{Critères d'acceptation :}

\begin{itemize}
\tightlist
\item
  Doit tenir compte des frais d'utilisation et d'une taxe à l'étranger.
\end{itemize}

\hypertarget{us10-activer-une-carte}{%
\subparagraph{US10 Activer une carte}\label{us10-activer-une-carte}}

En tant que titulaire possédant une carte, je veux pouvoir activer une
carte, afin d'accéder à toutes les fonctionalités possibles avec la
carte.

\textbf{Critères d'acceptation :}

\begin{itemize}
\tightlist
\item
  Activer une carte se fait directement dans l'application avec le code
  PIN.
\end{itemize}

\hypertarget{us11-modifier-le-code-pin-dune-carte}{%
\subparagraph{US11 Modifier le code PIN d'une
carte}\label{us11-modifier-le-code-pin-dune-carte}}

En tant que titulaire possédant une carte, Je veux pouvoir modifier le
code PIN de ma carte, afin de mémoriser plus facilement le code.

\textbf{Critères d'acceptations :}

\begin{itemize}
\tightlist
\item
  Cette opération requiert l'ancien code PIN.
\item
  Si la carte est inactive cela activera explicitement la carte avec une
  confirmation du titulaire.
\end{itemize}

\hypertarget{us12-validation-du-bloquage-dune-carte}{%
\subparagraph{US12 Validation du bloquage d'une
carte}\label{us12-validation-du-bloquage-dune-carte}}

En tant qu'employé d'une institution financiaire, je veux pouvoir
valider la demande de bloquage d'une carte d'un titulaire, afin de
rendre cette carte définitivement inactive.

\textbf{Critères d'acceptation :}

\begin{itemize}
\tightlist
\item
  Le titulaire de la carte bloquée reçoit une notifiaction lors de la
  validation du bloquaga.
\item
  Démarre la procédure de création d'une nouvelle carte si le client la
  demandée.
\end{itemize}

\hypertarget{us13-validation-dune-demande-de-carte}{%
\subparagraph{US13 Validation d'une demande de
carte}\label{us13-validation-dune-demande-de-carte}}

En tant qu'employé d'une institution financiaire, je veux pouvoir
valider la demande de création d'une nouvelle carte d'un client, afin
que le client reçoive une carte.

\textbf{Critères d'acceptation :}

\begin{itemize}
\tightlist
\item
  Une carte créée est inactive.
\end{itemize}

\hypertarget{us13-validation-dune-demande-de-modification-de-paramuxe8tres-dune-carte}{%
\subparagraph{US13 Validation d'une demande de modification de
paramètres d'une
carte}\label{us13-validation-dune-demande-de-modification-de-paramuxe8tres-dune-carte}}

En tant qu'employé d'une institution financiaire, je veux pouvoir
valider la demander de modification de paramètres de carte d'un client,
afin que le client ait les paramètres de son choix.

\textbf{Critères d'acceptation :}

\begin{itemize}
\tightlist
\item
  Si la demande est temporaire, aucune vérifiaction n'est requise lors
  de la date d'expiration du paramètre qui revient à sa valeur
  précédente.
\end{itemize}

\hypertarget{us14-visualiser-les-cartes-en-circulation}{%
\subparagraph{US14 Visualiser les cartes en
circulation}\label{us14-visualiser-les-cartes-en-circulation}}

En tant qu'employé d'une institution financiaire, je veux pouvoir
visualiser les cartes émisent par cette institution financiaire, afin de
savoir combien de cartes sont en circulation.

\textbf{Critères d'acceptation :}

\begin{itemize}
\tightlist
\item
  Seul les cartes émisent par cette institution sont visible.
\item
  Les infomations ne contiennent pas des information sensible telque le
  code PIN.
\end{itemize}
